\chapter{Consumption}\label{ch:b2:consumption}

Household consumption is about two thirds of GDP in the United States (see the second chart), and how it responds to income is one of the oldest
questions of macroeconomics. The chart shows, quarter by quarter for fifty years, the growth of American households'
real spending and of their real income after taxes.

\begin{openingchart}{Quarterly growth of real personal consumption expenditure and of real disposable personal income,
  United States, 1959Q2--2009Q3, per cent (log changes).}{statsmodels bundled data (public domain), macrodata; our
  computation}
\begin{tikzpicture}
\begin{axis}[omrecord, width=0.86\linewidth, height=0.42\linewidth, xmin=1959, xmax=2010, ymin=-4, ymax=6,
    xlabel={Year}, ylabel={Per cent a quarter}, xtick={1960,1970,1980,1990,2000,2010},
    xticklabel style={/pgf/number format/1000 sep=}, legend style={at={(0.5,-0.3)}, anchor=north,
    legend columns=2, /tikz/every even column/.append style={column sep=10pt}}]
  \addplot[omThm!70, thin] table[x=t, y=dy] {figdata/bachelor-2/19-consumption/out/growth.csv};
  \addplot[omDef, thick] table[x=t, y=dc] {figdata/bachelor-2/19-consumption/out/growth.csv};
  \legend{Disposable income, Consumption}
\end{axis}
\end{tikzpicture}
\end{openingchart}

% figdata: growth_summary (n 202, sd_dc 0.69, sd_dy 0.9, corr 0.44, mean_dc 0.84, mean_dy 0.83)
Both grew by about 0.8 per cent a quarter on average. But income growth varied more, with a standard deviation of 0.90
percentage points against 0.69 for consumption, and the two were only loosely related, with a correlation of 0.44.
Households did not spend each quarter what they earned that quarter; they smoothed.

The simplest consumption function, consumption as a fixed fraction of current income plus a constant, cannot produce
this. This chapter builds the models that can, from the household's choice over time of \cref{ch:b2:intertemporal-choice},
derives what they predict for the path of consumption, and tests them. Where they fail, the failure is informative: many
households spend much of a one-off payment as soon as it arrives.

\begin{recall}
From the school volumes: the consumption function and the marginal propensity to consume. From
\cref{ch:b2:intertemporal-choice}: the intertemporal budget constraint, the \omterm{def:b2:intertemporal-choice:euler-equation}{Euler equation} and \omterm{def:b2:intertemporal-choice:smoothing}{consumption smoothing}.
From \cref{ch:b2:expected-utility}: expected utility and the convexity of marginal utility.
\end{recall}

\begin{notation}[Symbols of this chapter]\label{not:b2:consumption:symbols}
$c_t$ is consumption and $y_t$ income in period $t$; $A_t$ financial wealth at the start of $t$; $r$ the real interest rate; $\beta$ the
\omterm{def:b2:intertemporal-choice:discount-factor}{discount factor}; $y^P$ \omterm{def:b2:consumption:permanent-income}{permanent income}; $\lambda$ the share of income going to households that spend their current income.
\end{notation}

\section{The simple consumption function against the data}\label{sec:b2:consumption:keynesian}

\begin{definition}[Average propensity to consume]\label{def:b2:consumption:apc}
The \emph{average propensity to consume}\index{average propensity to consume} is the ratio of consumption to income,
$c/y$.
\end{definition}

A consumption function $c = a + by$ with $a > 0$ and $0 < b < 1$ has an average propensity $a/y + b$ that falls as income rises toward $b$. As incomes
grow, a country's consumption should then take a falling share of its income. The chart shows household consumption as a
share of GDP in four countries over half a century in which real GDP per person rose 2.6 and 2.7 times in the United States
and Japan and 18.7 times in Korea (Germany, shown from 1991 only, 1.4 times since then).

\begin{omchart}{Household final consumption expenditure as a share of GDP, per cent, 1970 to the latest year; Germany from
  1991, since earlier figures are not data for the Germany of the time.}{World Bank,
  World Development Indicators (NE.CON.PRVT.ZS)}
\begin{tikzpicture}
\begin{axis}[omrecord, width=0.8\linewidth, height=0.42\linewidth, xmin=1970, xmax=2025, ymin=40, ymax=80,
    xlabel={Year}, ylabel={Per cent of GDP}, xtick={1970,1980,1990,2000,2010,2020},
    xticklabel style={/pgf/number format/1000 sep=}, legend style={at={(0.5,-0.3)}, anchor=north,
    legend columns=4, /tikz/every even column/.append style={column sep=8pt}}]
  \addplot[omDef, thick] table[x=year, y=USA] {figdata/bachelor-2/19-consumption/out/share.csv};
  \addplot[omThm, thick, dashed] table[x=year, y=JPN] {figdata/bachelor-2/19-consumption/out/share.csv};
  \addplot[omMeth, thick, dotted] table[x=year, y=KOR] {figdata/bachelor-2/19-consumption/out/share.csv};
  \addplot[black!60, thick, dashdotted] table[x=year, y=DEU] {figdata/bachelor-2/19-consumption/out/share.csv};
  \legend{United States, Japan, Korea, Germany}
\end{axis}
\end{tikzpicture}
\end{omchart}

% figdata: share_summary (USA 60.3 to 67.9, JPN 44.6 to 53.1, KOR 73.5 to 48, DEU 55.5 (1991) to 53.2; gdp_ratio 2.6, 2.7, 18.7, 1.4)
In Korea the share fell from 73.5 to 48 per cent; in the United States it rose from
60.3 to 67.9; in Germany, since 1991, it barely moved. The share is not a stable function of the level of income. What households spend
depends on more than this year's income.

\section{The life-cycle and permanent-income models}\label{sec:b2:consumption:pih}

\begin{definition}[Life-cycle hypothesis]\label{def:b2:consumption:life-cycle}
The \emph{life-cycle hypothesis}\index{life-cycle hypothesis} states that households plan consumption over their whole
lives, saving while they earn and drawing on their savings in retirement, so that consumption follows lifetime
resources rather than current income.
\end{definition}

\begin{definition}[Permanent income]\label{def:b2:consumption:permanent-income}
A household's \emph{permanent income}\index{permanent income} is the constant income whose present value over its
remaining life equals its wealth plus the present value of its expected income.
\end{definition}

\begin{definition}[Permanent-income hypothesis]\label{def:b2:consumption:pih}
The \emph{permanent-income hypothesis}\index{permanent-income hypothesis} states that households consume their \omterm{def:b2:consumption:permanent-income}{permanent
income}, so that consumption responds fully to changes in \omterm{def:b2:consumption:permanent-income}{permanent income} and little to changes in income that are
expected to be temporary.
\end{definition}

Both follow from the household's problem of \cref{ch:b2:intertemporal-choice}. If $\beta(1 + r) = 1$, the \omterm{def:b2:intertemporal-choice:euler-equation}{Euler equation} $u'(c_t) =
\beta(1 + r)u'(c_{t+1})$ gives $c_t = c_{t+1}$: consumption is flat, and the budget constraint fixes its level. With $r = 0$ for simplicity and
$T$ years left, $c = (A_0 + \sum_t y_t)/T$.

Take a worker who will earn 30 thousand dalens a year for 40 years and then live 20 years in retirement, with no wealth to
start. She consumes $40 \times 30/60 = 20$ a year, saves 10 a year while working, accumulates 400 by retirement and runs it down to
zero.

\begin{omfigure}
\begin{tikzpicture}[econ]
\begin{axis}[omrecord, grid=none, xmin=0, xmax=62, ymin=0, ymax=36, width=0.62\linewidth,
    height=0.4\linewidth, xlabel={Years from the start of working life}, ylabel={Thousands of dalens a year},
    xtick={0,40,60}, ytick={20,30}, legend style={at={(1.03,0.5)}, anchor=west}]
  \addplot[omThm, very thick, const plot] coordinates {(0,30) (40,0) (60,0)};
  \addplot[omDef, very thick, dashed, domain=0:60, samples=2] {20};
  \addplot[fill=omDef!12, draw=none, forget plot] coordinates {(0,20) (40,20) (40,30) (0,30)} -- cycle;
  \addplot[fill=omThm!12, draw=none, forget plot] coordinates {(40,0) (60,0) (60,20) (40,20)} -- cycle;
  \node[font=\footnotesize] at (axis cs:20,25) {saving};
  \node[font=\footnotesize] at (axis cs:50,10) {dissaving};
  \legend{Income, Consumption}
\end{axis}
\end{tikzpicture}

{\small The life cycle. Income is 30 for forty years and zero in retirement; consumption is flat at 20. The household saves
10 a year while working and draws down its savings, 400 at retirement, over the last twenty years.}
\end{omfigure}

The models imply that the marginal propensity to consume depends on how long a change in income lasts. A one-off payment
of $W$ raises consumption by $W/T$ a year: with fifty years left, by 2 per cent of the payment. A raise that lasts the rest of the
working life raises it by the share of remaining life spent working (half of it in the chart below, two thirds in the
example above), and a raise for life by all of it.

\begin{omfigure}
\begin{tikzpicture}[econ]
\begin{axis}[omrecord, ybar, bar width=18pt, width=0.55\linewidth, height=0.4\linewidth, ymin=0, ymax=1.15,
    xmin=0.4, xmax=3.6, xtick={1,2,3}, xticklabels={One-off, Working life, For life},
    ylabel={Rise in consumption per dalen of income}, xmajorgrids=false, nodes near coords,
    nodes near coords style={font=\footnotesize, /pgf/number format/fixed, /pgf/number format/precision=2}]
  \addplot[fill=omDef!50, draw=omDef] coordinates {(1,0.02) (2,0.5) (3,1)};
\end{axis}
\end{tikzpicture}

{\small The marginal propensity to consume out of a change in income, for a household with fifty years left of which
twenty-five are working years, $r = 0$: a one-off payment, a raise for the rest of working life, and a raise for life
(including the pension).}
\end{omfigure}

\begin{model}{Life-cycle and permanent income}
\assumes Households that plan over their remaining lives, borrow and lend at one interest rate, and know the income they
can expect.
\predicts Consumption smoother than income; saving in working life and dissaving in retirement; a small response to
temporary income and a full one to \omterm{def:b2:consumption:permanent-income}{permanent income}; no response when an anticipated payment arrives.
\tested Consumption growth is smoother than income growth (the opening chart); lagged income does not predict consumption
growth in postwar American data (Hall 1978).
\breaks When one-off payments arrive: households spend a large share of tax rebates within months (Johnson, Parker and
Souleles 2006).
\end{model}

\section{The Euler equation and the random walk}\label{sec:b2:consumption:random-walk}

With uncertain income the \omterm{def:b2:intertemporal-choice:euler-equation}{Euler equation} holds in expectation: $u'(c_t) = \beta(1 + r)E_t u'(c_{t+1})$. If $\beta(1 + r) = 1$ and utility is quadratic, so
that marginal utility is linear in consumption, $c_t = E_t c_{t+1}$.

\begin{definition}[Random-walk hypothesis of consumption]\label{def:b2:consumption:random-walk-hypothesis}
The \emph{random-walk hypothesis}\index{random-walk hypothesis} of consumption states that the best forecast of next
period's consumption is this period's consumption, so that changes in consumption cannot be predicted from anything
known today.
\end{definition}

\begin{proposition}[Consumption is a random walk]\label{prop:b2:consumption:random-walk}
If utility is quadratic, $\beta(1 + r) = 1$ and the household can borrow and lend freely at $r$, then $c_{t+1} = c_t + \varepsilon_{t+1}$, where $\varepsilon_{t+1}$ has mean zero given
all information at $t$. In particular, predictable changes in income do not predict changes in consumption.
\end{proposition}

\begin{proof}
With $u(c) = c - \tfrac{a}{2}c^2$, $u'(c) = 1 - ac$, and the \omterm{def:b2:intertemporal-choice:euler-equation}{Euler equation} with $\beta(1 + r) = 1$ gives $1 - ac_t = 1 - aE_t c_{t+1}$, so $E_t c_{t+1} = c_t$. Let $\varepsilon_{t+1} = c_{t+1} - E_t c_{t+1}$;
its expectation at $t$ is zero by construction, so it is uncorrelated with anything known at $t$, including the expected change in
income. A predictable rise in income was already in the household's \omterm{def:b2:consumption:permanent-income}{permanent income}, and so already in $c_t$.
\end{proof}

\begin{definition}[Excess sensitivity]\label{def:b2:consumption:excess-sensitivity}
\emph{Excess sensitivity}\index{excess sensitivity} of consumption is a response of consumption to changes in income
that were predictable in advance, which the \omterm{def:b2:consumption:random-walk-hypothesis}{random-walk hypothesis} rules out.
\end{definition}

The test regresses consumption growth on the growth of income that was predictable: if some households spend their
current income, a share $\lambda$ of income growth passes into consumption growth, $\Delta c_t = \mu + \lambda\Delta y_t + \varepsilon_t$. Because income growth is
correlated with the news in $\varepsilon_t$, ordinary least squares is biased; the test uses only the part of income growth that could
be predicted from variables known two or more quarters earlier, an instrumental-variable estimate. The listing runs it
on the data of the opening chart.

\omcode{code/bachelor-2/19-consumption/listings/lambda_iv.py}{1}{33}{Estimating the share of rule-of-thumb consumers
by instrumental variables on quarterly United States data.}\label{lst:b2:consumption:lambda}

\omoutput{code/bachelor-2/19-consumption/outputs/lambda_iv.txt}{The instrumented share is about 0.6, but the
instruments predict income growth poorly: the first-stage $F$ is 2.5.}\label{out:b2:consumption:lambda}

% code: listings/lambda_iv.py
% ledger: b2:cm1989, b2:hall1978
The estimate, 0.62 with a standard error of 0.20, rejects $\lambda = 0$, the random walk. But the six instruments explain only 7 per
cent of income growth, with an $F$ statistic of 2.5, far below the level at which instrumental-variable estimates are
reliable; such weak instruments pull the estimate toward the least-squares slope and make its standard error
misleading. Campbell and Mankiw (1989) read the American evidence as about half
of consumers spending their current income. Hall (1978) had found that lagged income did not predict consumption, but
stock prices lagged by one quarter did.

\section{Precautionary saving and borrowing limits}\label{sec:b2:consumption:precautionary}

Quadratic utility is convenient but implies that uncertainty does not change saving. With marginal utility that is
convex, as for any utility with \omterm{def:b2:expected-utility:crra}{constant relative risk aversion}, it does.

\begin{definition}[Prudence]\label{def:b2:consumption:prudence}
A household is \emph{prudent}\index{prudence} if its marginal utility is convex, $u''' > 0$, so that a spread in future
consumption raises expected marginal utility.
\end{definition}

\begin{definition}[Precautionary saving]\label{def:b2:consumption:precautionary}
\emph{Precautionary saving}\index{precautionary saving} is the extra saving caused by uncertainty about future income,
over what a household would save if the same expected income were certain.
\end{definition}

Take two periods, $r = 0$, $\beta = 1$ and $u = \ln c$. Income is 100 now and 100 next period. With certainty the household saves nothing. If next period's income
is 50 or 150, equally likely, the \omterm{def:b2:intertemporal-choice:euler-equation}{Euler equation} $1/(100 - s) = \tfrac12/(50 + s) + \tfrac12/(150 + s)$ reduces to $s^2 + 100s - 1{,}250 = 0$, so $s = 11.24$. The household saves
11 per cent of its income though its expected future income has not changed, because a bad outcome raises marginal utility
more than a good one lowers it: $\tfrac12(1/50 + 1/150) = 0.0133$, against $1/100 = 0.01$.

\begin{omfigure}
\begin{tikzpicture}[econ]
\begin{axis}[omrecord, grid=none, xmin=30, xmax=170, ymin=0, ymax=0.026, width=0.6\linewidth,
    height=0.4\linewidth, xlabel={Consumption}, ylabel={Marginal utility $1/c$}, xtick={50,100,150},
    ytick={0.0067,0.01,0.0133,0.02}, yticklabels={0.0067,0.0100,0.0133,0.0200},
    scaled y ticks=false]
  \addplot[omDef, very thick, domain=40:170, samples=80] {1/x};
  \draw[black!60] (axis cs:50,0.02) -- (axis cs:150,0.006667);
  \node[eqpoint] at (axis cs:100,0.01) {};
  \node[eqpoint] at (axis cs:100,0.013333) {};
  \draw[guide] (axis cs:100,0) -- (axis cs:100,0.013333);
  \node[anchor=west, font=\footnotesize] at (axis cs:102,0.0138) {expected marginal utility};
  \node[anchor=north east, font=\footnotesize] at (axis cs:99,0.0098) {at the mean};
\end{axis}
\end{tikzpicture}

{\small Why uncertainty raises saving. With convex marginal utility, the average of marginal utility at 50 and 150 (on the
chord) is above marginal utility at their mean, 100. To restore the \omterm{def:b2:intertemporal-choice:euler-equation}{Euler equation}, the household moves resources to the
future: it saves.}
\end{omfigure}

A borrowing limit has a related effect. A household whose income will rise but who cannot borrow against it consumes its
current income, and spends any extra cash as it arrives, until it reaches the consumption it would choose if it could
borrow. Together, \omterm{def:b2:consumption:prudence}{prudence} and borrowing limits make households hold a small stock of savings as a buffer against bad
times, and make those with little of it spend windfalls quickly.

\section{Where it breaks: one-off payments}\label{sec:b2:consumption:rebates}

\begin{definition}[Hand-to-mouth household]\label{def:b2:consumption:hand-to-mouth}
A \emph{hand-to-mouth household}\index{hand-to-mouth household} is one that consumes all or nearly all of its current
disposable income in each period, because it has little liquid wealth or cannot borrow.
\end{definition}

\begin{record}[The 2001 and 2008 tax rebates]
% ledger: b2:jps2006, b2:psjm2013
In 2001 and again in 2008 the United States government sent households one-off payments, and the date each household
received its payment was set in a way unrelated to its circumstances, so that households paid earlier can be compared
with households paid later. In 2001 households spent 20 to 40 per cent of their rebates on nondurable goods in the three
months in which they arrived, and about two thirds over six months; responses were larger for households with little
liquid wealth or low income (Johnson, Parker and Souleles 2006). In 2008 households spent 12 to 30 per cent of their payments on
nondurable goods in the three months of receipt, and with durable goods, mostly vehicles, 50 to 90 per cent in all
(Parker, Souleles, Johnson and McClelland 2013). Under the permanent-income model, a one-off payment should be spread over
the rest of a household's life, so that a few per cent at most is spent in the first quarter.
\end{record}

If a share $h$ of households are hand-to-mouth and spend a one-off payment in full, and the rest follow the permanent-income
model, the aggregate propensity to consume a payment within a quarter is $h + (1 - h)m$, where $m$ is the forward-looking households'
quarterly propensity. With fifty years left and $r = 0$, $m = 1/200 = 0.005$; a response of 20 to 40 per cent then implies that 20 to 40 per cent of
households live hand to mouth.

\begin{measure}[The household saving rate]
The household saving rate is disposable income minus consumption, divided by disposable income. Both terms are measured in
the national accounts, which estimate them from surveys of spending, tax records and business data, so the saving rate is
a small difference between two large estimates: an error of one per cent in consumption moves a saving rate of 5 per cent
by almost one point. The accounts also count as consumption the purchase of durable goods, such as cars, that serve for
years, though buying one is in part a way of carrying resources into the future.
\end{measure}

\section{Exercises}\label{sec:b2:consumption:exercises}

\begin{exercise}[$\star$]\label{exo:b2:consumption:1}
Read the opening chart and its summary. Which series is smoother, and by how much? What would a consumption function
$c = 0.9y$ predict for the two standard deviations and their correlation?
\end{exercise}

\begin{exercise}[$\star$]\label{exo:b2:consumption:2}
A worker will earn 36 thousand dalens a year for 30 years and then live 15 years in retirement; she has wealth 90; $r = 0$. Find her
consumption, her saving while working and her wealth at retirement.
\end{exercise}

\begin{exercise}[$\star$]\label{exo:b2:consumption:3}
Read the record on rebates. Why was the date of receipt being unrelated to circumstances essential to the studies?
\end{exercise}

\begin{exercise}[$\star$]\label{exo:b2:consumption:4}
Predict, then compute: if the quarterly response to a payment is 0.12 and forward-looking households have $m = 0.005$, what share of
households are hand to mouth?
\end{exercise}

\begin{exercise}[$\star\star$]\label{exo:b2:consumption:5}
In the precautionary example, let next period's income be 25 or 175. Find the saving and compare it with the example.
\end{exercise}

\begin{exercise}[$\star\star$]\label{exo:b2:consumption:6}
A household lives two periods with incomes 20 and 40, $r = 0$, and cannot borrow. Find its consumption now, and its response to a
one-off payment now of 6, of 10 and of 30.
\end{exercise}

\begin{exercise}[$\star\star$]\label{exo:b2:consumption:7}
Explain why ordinary least squares overstates or understates $\lambda$ in the regression of consumption growth on income growth,
and why weak instruments are a problem. What does the listing's first-stage $F$ say?
\end{exercise}

\begin{exercise}[$\star\star$]\label{exo:b2:consumption:8}
Modify the listing to use only lagged income growth as instruments. Report $\lambda$ and the first-stage $F$, and explain the change.
\end{exercise}

\begin{exercise}[$\star\star$]\label{exo:b2:consumption:9}
Read the chart of consumption shares. Give two reasons, other than the level of income, why the share might fall in one
country and rise in another.
\end{exercise}

\begin{exercise}[$\star\star\star$]\label{exo:b2:consumption:10}
With $r > 0$, $\beta(1 + r) = 1$ and an infinite horizon, show that consumption equals $rA_t/(1 + r)$ plus the annuity value of expected income, and
that a one-off payment $W$ raises consumption by $rW/(1 + r)$ a year.
\end{exercise}

\begin{exercise}[$\star\star\star$]\label{exo:b2:consumption:11}
Show that with $u = \ln c$, $r = 0$ and $\beta = 1$, the saving of the two-period example solves $s^2 + 100s - 1{,}250 = 0$ when income next period is 50 or 150, and that
saving rises with the spread when income is $100 - \delta$ or $100 + \delta$.
\end{exercise}

\begin{exercise}[$\star\star\star$]\label{exo:b2:consumption:12}
Show that if a share $\lambda$ of income goes to households that consume their current income and the rest to households whose
consumption is a random walk, then $\Delta c_t = \lambda\Delta y_t + (1 - \lambda)\varepsilon_t$, and that $\lambda$ is the coefficient on predictable income growth.
\end{exercise}

\section{Problem: A Brenland household gets a bonus}\label{sec:b2:consumption:problem}

\begin{problem}[{Weekend problem --- a Brenland household receives a bonus: how much does it spend under the permanent-income
model, with a borrowing limit, and across a population in which some households live hand to mouth}]
\label{pb:b2:consumption:1}
Amounts are in thousands of dalens; the interest rate is zero and the household does not discount.

\medskip
\noindent\textbf{Part I --- Permanent income.} A household has 50 years to live, of which 25 are working years at an income of 30
a year; it has wealth 250.
\begin{enumerate}
  \item What are its lifetime resources?
  \item What does it consume each year?
  \item How much does it save each working year?
  \item What is its wealth at retirement?
  \item How does its wealth evolve in retirement?
\end{enumerate}

\medskip
\noindent\textbf{Part II --- Transitory and permanent changes.}
\begin{enumerate}[resume]
  \item It receives a one-off bonus of 10. By how much does consumption rise, and what is the propensity to consume?
  \item Instead it receives a raise of 2 a year for its 25 working years. By how much does consumption rise?
  \item Instead the raise also raises its retirement income by 2 a year. By how much does consumption rise?
  \item The bonus of 10 is announced a year before it is paid. When does consumption rise?
  \item What would \omterm{def:b2:consumption:excess-sensitivity}{excess sensitivity} look like here?
  \item Why does the propensity in Question 6 overstate the response within a quarter?
\end{enumerate}

\medskip
\noindent\textbf{Part III --- A borrowing limit.} A young household lives two periods with incomes 20 and 40 and cannot
borrow.
\begin{enumerate}[resume]
  \item What would it consume each period if it could borrow?
  \item What does it consume now?
  \item How much does it spend of a bonus of 6 received now?
  \item Of a bonus of 10?
  \item Of a bonus of 30?
  \item Why does its propensity to spend fall as the bonus grows?
\end{enumerate}

\medskip
\noindent\textbf{Part IV --- A population.} Brenland sends a one-off payment of 0.6 (600 dalens) to each of 2{,}000 households. Forward-looking
households are like the household of Part I; a share $h$ live hand to mouth.
\begin{enumerate}[resume]
  \item What is the forward-looking households' propensity to spend the payment within a quarter?
  \item If total spending in the quarter is 30 per cent of the payments, how much is spent?
  \item What share $h$ does that imply?
  \item What range of $h$ does a quarterly response of 20 to 40 per cent imply?
  \item If $h$ were one half, what would the aggregate propensity be?
  \item Why might households with large but illiquid wealth behave as hand to mouth?
  \item Why does $h$ matter for a government deciding whether to send such payments in a recession?
  \item State the result: the aggregate propensity to consume a one-off payment and the share of \omterm{def:b2:consumption:hand-to-mouth}{hand-to-mouth households} it
    implies.
\end{enumerate}
\end{problem}
